\section{Combiner in the Deep Kalman Filter (DKF)}\label{combiner_info}
    
The \textbf{Combiner} is a key component in the DKF model, responsible 
for integrating information from different sources to produce a unified 
representation. Specifically, it combines the previous latent state 
($\bm{z}_{t-1}$) with the hidden states of one or more Recurrent Neural 
Networks (RNNs) to create an output that is fed into the encoder. 
This output is then used to compute the parameters of the latent 
distribution ($\bm{z}_t$).

\subsection*{Purpose of the Combiner}

The Combiner's primary role is to:
\begin{enumerate}
    \item \textbf{Integrate Information}: Combine the previous latent 
    state ($\bm{z}_{t-1}$) with the hidden states of RNNs (e.g., backward RNN, 
    forward RNN) to create a richer representation.
    \item \textbf{Produce a Unified Output}: Generate a single tensor that 
    can be fed into the encoder to compute the mean and log variance of the 
    latent state ($\bm{z}_t$).
    \item \textbf{Flexibility}: Allow for different ways of combining 
    information, such as linear combination or using a Multi-Layer 
    Perceptron (MLP).
\end{enumerate}

\subsection*{Key Concepts}

\begin{itemize}
    \item \textbf{Latent State ($\bm{z}_{t-1}$)}:
        \begin{itemize}
            \item The latent state from the previous time step, which 
            represents the hidden variables of the model.\\
            At $t=1$, $\bm{z}_0$ is either a fix tensor of zeros or a learnable
            parameter initialiazed with zeros.
        \end{itemize}
    \item \textbf{RNN Hidden States ($\bm{h}_{t}^{i}$)}:
        \begin{itemize}
            \item The output of one or more RNNs (e.g., backward RNN, 
            forward RNN) that process the input sequence. These hidden 
            states capture temporal dependencies in the data.
        \end{itemize}
    \item \textbf{Combination Strategy}:
        \begin{itemize}
            \item The method used to combine $\bm{z}_{t-1}$ with the RNN 
            hidden states. This can be a simple linear combination or 
            a more complex MLP-based combination.
        \end{itemize}
\end{itemize}

\subsection*{Combiner Types}

The Combiner is implemented in two main ways:
\begin{enumerate}
    \item \textbf{Linear Combiner}: Combines $\bm{z}_{t-1}$ and RNN hidden 
    states using a weighted sum (linear combination).
    \item \textbf{MLP Combiner}: Combines $\bm{z}_{t-1}$ and RNN hidden 
    states using a Multi-Layer Perceptron (MLP).
\end{enumerate}

\subsection{Linear Combiner}

The \textbf{Linear Combiner} performs a weighted sum of the previous 
latent state ($\bm{z}_{t-1}$) and the RNN hidden states. 
The weights are controlled by a parameter $\alpha_i$, which determines 
the contribution of each input.

\subsubsection*{Key Parameters}
\begin{itemize}
    \item \code{$z\_dim$}: Dimension of the latent state ($\bm{z}_{t-1}$).
    \item \code{$rnn\_dims$}: Dimensions of the RNN hidden states. 
    This can be a single dimension or a list of dimensions if multiple 
    RNNs are used.
    \item \code{$output\_dim$}: Dimension of the output tensor.
    \item \code{$\alpha^i$}: Weight for the $i$-th RNN hidden state. 
    This determines how much influence the mapped $\bm{z}_{t-1}$ has on 
    the combination.
\end{itemize}

\subsubsection*{Formula}

\begin{equation*}
    \bm{combined}^{i}_{t} = \alpha^{i} \cdot \text{mapped}(\bm{z}_{t-1}) + 
    (1 - \alpha^{i}) \cdot \bm{h}_{t}^{i}
\end{equation*}

Where:

\begin{itemize}
    \item $\text{mapped}(\bm{z}_{t-1}) = \tanh(\mathbf{W}^i \bm{z}_{t-1}
    + \mathbf{c}^i)$: the previous latent state mapped to the dimension of
    the $i$-th RNN hidden state by a linear layer followed by the
    activation function ($\tanh$ in all experiments).
    \item $\bm{h}_{t}^{i}$: The hidden state output of the $i$-th RNN
    at time $t$.
\end{itemize}

With a single RNN and $\alpha = 0.5$, this is the combiner of Krishnan
et al.'s follow-up work \cite{krishnan2017structured},
$\frac{1}{2}\big(\tanh(\mathbf{W}\bm{z}_{t-1} + \mathbf{c}) +
\bm{h}_t\big)$, followed by a linear layer.

Two sets of weights are kept: \code{alpha\_init}, meant for the first
time step, where $\bm{z}_{0}$ is a fixed or learned constant that carries
no information about the sequence, and \code{alpha}, used for the
following steps.

\subsubsection*{How It Works}
\begin{enumerate}
    \item \textbf{Mapping $\bm{z}_{t-1}$}:
        \begin{itemize}
            \item The previous latent state ($\bm{z}_{t-1}$) is mapped to
            the dimension of each RNN hidden state by separate linear
            layers, each followed by the activation function.
        \end{itemize}
    \item \textbf{Weighted Sum}:
        \begin{itemize}
            \item The mapped $\bm{z}_{t-1}$ and the RNN hidden state are 
            combined using a weighted sum. The weight $\alpha^i$ 
            controls the contribution of the mapped $\bm{z}_{t-1}$, while 
            $(1 - \alpha^i)$ controls the contribution of the RNN hidden 
            state.
        \end{itemize}
    \item \textbf{Concatenation}:
        \begin{itemize}
            \item The combined tensors for all RNN hidden states are 
            concatenated along the feature dimension.
        \end{itemize}
    \item \textbf{Final Linear Layer}:
        \begin{itemize}
            \item The concatenated tensor is passed through a final linear
            layer to produce the output tensor of dimension \code{output\_dim}.
        \end{itemize}
\end{enumerate}

\subsubsection*{Example}
Suppose:

\begin{itemize}
    \item $\alpha^i = 0.3$
    \item $\text{mapped}(\bm{z}_{t-1}) = \bm{mz}$
    \item $\bm{h}_{t}^{i} = \bm{rnn\_h}$
\end{itemize}

The formula becomes:

\begin{align*}
    \bm{combined}^i &= 0.3 \cdot \bm{mz} + 
    (1 - 0.3) \cdot \bm{rnn\_h} \\
    &= 0.3 \cdot \bm{mz} + 0.7 \cdot \bm{rnn\_h}
\end{align*}

This means:
\begin{itemize}
    \item 30\% of the contribution comes from the mapped latent state 
    $\text{mapped}$.
    \item 70\% of the contribution comes from the RNN hidden state 
    \code{rnn\_h}.
\end{itemize}

\subsection{MLP Combiner}

The \textbf{MLP Combiner} uses a Multi-Layer Perceptron (MLP) to combine
the previous latent state ($\bm{z}_{t-1}$) and the RNN hidden states.
This allows for more complex interactions between the inputs.

\subsubsection*{Key Parameters}
\begin{itemize}
    \item \code{z\_dim}: Dimension of the latent state ($\bm{z}_{t-1}$).
    \item \code{rnn\_dims}: Dimensions of the RNN hidden states.
    \item \code{output\_dim}: Dimension of the output tensor.
    \item \code{layers\_dim}: List of dimensions for the hidden
    layers in the MLP.
    \item \code{activation}: Activation function 
    (e.g., $\tanh$ or $\text{ReLU}$).
\end{itemize}

\subsubsection*{How It Works}
\begin{enumerate}
    \item \textbf{Concatenation}:
        \begin{itemize}
            \item The previous latent state ($\bm{z}_{t-1}$) and all RNN hidden
            states are concatenated along the feature dimension.
        \end{itemize}
    \item \textbf{MLP Layers}:
        \begin{itemize}
            \item The concatenated tensor is passed through a series of
            linear layers with activation functions in between.
            \item The first layer maps the concatenated input to the first
            hidden layer dimension.
            \item Intermediate layers map the output of one hidden layer to
            the next.
            \item The final layer maps the last hidden layer to the output
            dimension (\code{output\_dim}).
        \end{itemize}
    \item \textbf{Output}:
        \begin{itemize}
            \item The output of the MLP is the combined tensor, which is
            then fed into the encoder.
        \end{itemize}
\end{enumerate}

\subsubsection{Example}
Suppose:
\begin{itemize}
    \item $\bm{z}_{t-1}$ has a dimension of 16.
    \item Two RNN hidden states: $h^1_t$ (dimension 32) and $h^2_t$ (dimension 64).
\end{itemize}
%
The concatenated input will have a dimension of \(16 + 32 + 64 = 112\).
This concatenated tensor is passed through an MLP.

\subsection{Combiner Factory}

The \textbf{Combiner Factory} is a utility class that creates the
appropriate Combiner based on the specified type ($\text{linear}$ or
$\text{mlp}$). This allows for flexibility in choosing the combination
strategy.

\subsubsection*{Key Parameters}

\begin{itemize}
    \item \code{combiner\_type}: Type of Combiner to use 
    (\text{linear} or \text{mlp}).
    \item \textbf{Other Parameters}: Depend on the type of Combiner 
    (e.g., \code{$\alpha^i$}, \code{layers\_dim}).
\end{itemize}

\subsubsection*{How It Works}

\begin{itemize}
    \item If \code{combiner\_type} = \code{linear}, the factory 
    initializes a \code{DKF\_CombinerLinear} with the specified parameters.
    \item If \code{combiner\_type} = \code{mlp}, the factory initializes 
    a \code{DKF\_CombinerMLP} with the specified parameters.
\end{itemize}
\subsection{Summary of the Combiner's Role in the DKF}

The Combiner is responsible for integrating the previous latent state
($\bm{z}_{t-1}$) with the hidden states of RNNs to produce a unified
representation. This representation is then used by the encoder to
compute the parameters of the latent distribution ($\bm{z}_{t}$).
The Combiner can be either a simple linear combination or a more complex
MLP-based combination, depending on the model's requirements.
\\
By using the Combiner, the DKF model can effectively leverage both the
temporal dependencies captured by RNNs and the latent state dynamics to
model complex time series data.
